\scanpage{5}

% Two parallel pen strokes stand in the upper left margin.
\uline{Th} \quad $f:X\longrightarrow S$ morphisme \inserted{\uline{plat}} de type fini, avec
$S$ loc noeth, $Y$ \deleted{\unsure{sous-préschéma} \unsure{fermé} de $X$}\inserted{$\longrightarrow X$ un \unsure{mono|morph.} de type fini}.

Pour que $\displaystyle\prod_{X/S}Y/X$ soit représentable, il f. et s. qu'il commute aux limites projectives adiques d'anneaux artiniens\inserted{, si $Y\to X$ est une \unsure{imm} (\unsure{imm} \uuline{f.})}, et alors il est représentable par un sous-préschéma (\uuline{fermé}) de $S$.
% slanted
\margin{il suffit que $Y\to X$ soit un \uuline{monomorphisme}}
\centerline{\rule{0.94\linewidth}{0.3pt}}
\deleted{\unsure{fpqc|qui}}

\uline{Dém}. Vérifions les conditions générales du th. de représentabilité par morphismes non ramifiés séparés

$1^\circ)$ $F$ faisceau pour fpqc : \deleted{\unsure{EGA IV 2|EGA IV 8}} clair EGA IV 2

$2^\circ)$ $F$ commute aux $\underrightarrow{\lim}$ d'anneaux : clair EGA IV 8
% slanted
\margin{$X$ t.f.}

$3^\circ)$ C'est la condition de commutation aux $\underleftarrow{\lim}$ adiques d'artiniens.

$4^\circ)$, $\deleted{\unsure{6}}5^\circ)$, $6^\circ)$ Dans le cas d'un monomorphisme $F\longrightarrow S$ se réduit à ceci

si $A\longrightarrow A'$ est un hom injectif d'artiniens \inserted{$A'/A$ fini}
alors $F(A)\to F(A')$ est bijectif, ou encore $F(A')\ne\emptyset\Longrightarrow F(A)\unsure{=|\ne}\emptyset$.
Il faut prouver que si $X_{A'}=Y_{A'}$, on a $X_A=Y_A$. C'est clair, car $X_A$ \deleted{\unsure{étant}} plat sur $A$,
$\mathcal O_{X_A}\deleted{\unsure{\otimes_A}}\to\mathcal O_{X_A}\otimes_A A'$ est injectif.
% slanted
\margin{$X$ \uline{plat}}

$7^\circ)$ Si $T$ sur $S$ est \deleted{\unsure{\illegible{1}|noeth.}} local, noeth, muni d'un idéal de carré nul $J$, $T_0=V(J)$, $U$ ouvert de $T$ tel que $T$ en soit l'adhérence schématique, \uline{alors} pour
$\deleted{\unsure{\xi\in F|\xi\in\mathcal F|F(T)}\,\unsure{tF}}$
% slanted
\margin{$X$ \uline{plat} \deleted{\unsure{plat}}}
\diagram{\begin{tikzcd}[ampersand replacement=\&]
F(T)\arrow[r,dash]\arrow[d] \& F(T_0)\arrow[d]\\
F(U)\arrow[r] \& F(U_0)
\end{tikzcd}}
est cartésien.
% archive sheet number: 4
